\section{A path obstruction from the Horak--Kim construction}
\label{sec:second-order-constructions}

The proposed additive path bound needs a polynomial factor.
The following family gives a linear obstruction at terminal distance three.
It uses the graphs $MC_{2q}$ of Horak and Kim~\cite{HK}.
Their total cycle count already gives the known lower bound
$\psi(r)-2^{r-1}=\Omega(r2^{r/2})$.
The marked-edge count below gives a separate consequence for terminal paths.

\begin{proposition}\label{prop:hk-short-path-obstruction}
For every integer $q\ge2$, there is a simple connected graph $N_q$
with terminals $u,v$, rank $b=2q$, and $\operatorname{dist}_{N_q}(u,v)=3$,
such that
\[
 P_{uv}(N_q)=2^{2q-1}+\left(q-\frac34\right)2^q+1.
\]
Consequently, no absolute constant $C$ satisfies
\[
 P_{uv}(N)\le 2^{b-1}+C2^{(b+\ell)/2}
\]
for every connected graph of rank $b$ with terminal distance $\ell$.
\end{proposition}

\begin{proof}
Use indices modulo $q$.
Block $i$ has vertices $a_i,b_i,c_i,d_i$ and all four edges between
$\{a_i,b_i\}$ and $\{c_i,d_i\}$.
Add the interface edges $c_i a_{i+1}$ and $d_i b_{i+1}$.
The resulting graph $G_q$ is simple and cubic.
It has $4q$ vertices, $6q$ edges, and rank $2q+1$.
Set $e=a_0c_0$ and $N_q=G_q-e$.

Cut parity gives two cases for every cycle of $G_q$.
It uses one edge at every interface, or it uses zero or two at every interface.
In the first case, choose one edge independently at every interface.
Each square then has two paths between the selected endpoints.
There are therefore $2^q2^q=4^q$ cycles in this case.
Exactly half contain $e$.
Indeed, four of the eight local endpoint-and-path choices contain any fixed square edge.
The contribution is $2^{2q-1}$.

Now consider the second case.
If no interface is used, only the square in block zero contains $e$.
This contributes one cycle.
If some interfaces are used, but not all, they form a cyclic interval.
An interval of $d$ interfaces occupies $d+1$ blocks, where $1\le d\le q-1$.
Each end block has two possible caps.
Each inner block has two possible matchings.
There are $2^{d+1}$ cycles for each such interval.
Exactly $d+1$ intervals contain block zero.
For each of these intervals, half the cycles contain $e$.
Their total contribution is
\[
 \sum_{d=1}^{q-1}(d+1)2^d=(q-1)2^q.
\]

Finally, suppose that every interface is used.
Each square supplies one of its two matchings.
The result is one cycle precisely when the matching product is the transposition.
There are $2^{q-1}$ such choices.
Exactly half contain $e$, so they contribute $2^{q-2}$.
Adding the contributions gives
\[
 c_e(G_q)=2^{2q-1}+1+(q-1)2^q+2^{q-2}.
\]

Deleting $e$ gives a bijection from these cycles to simple terminal paths.
The remaining three edges of its square give a path of length three.
The graph is bipartite, so the terminal distance is three.
Its rank is $2q$.
The normalized excess equals
\[
 \frac{P_{uv}(N_q)-2^{b-1}}{2^{(b+3)/2}}
 =\frac{q-3/4+2^{-q}}{2^{3/2}},
\]
which tends to infinity.
\end{proof}

This obstruction does not refute an error bound $O(r2^{r/2})$ for total cycles.
It shows that a direct path bound at the half exponent needs additional rank dependence.
The independent script \texttt{check\_second\_order\_constructions.py}
counts terminal paths for $2\le q\le9$ and verifies the formula.
These finite checks support the symbolic proof; they do not establish optimality.
